% chktex-file 1
% chktex-file 3
% chktex-file 8
% chktex-file 36
% Suppress chktex false positives:
%   1  -- \skewt uses \xspace so trailing space is intentional
%   3  -- math like \{Y_t\}_{t=1}^{n_t} parses correctly
%   8  -- en-dashes for proper-name pairs and ranges; ISO dates with hyphen
%   36 -- AR(p) is canonical notation; no space before the parenthesis
\documentclass[11pt]{article}

\usepackage[margin=1in]{geometry}
\usepackage{amsmath,amssymb,amsthm}
\usepackage{bm}
\usepackage{xspace}
\usepackage{booktabs}
\usepackage{enumitem}
\usepackage{listings}
\usepackage{xcolor}
\usepackage{hyperref}
\usepackage[capitalize]{cleveref}
\hypersetup{colorlinks=true, linkcolor=blue, urlcolor=blue, citecolor=blue}

% ---- Notation aligned with skewt_rev_2.tex / mycommands.sty ----
\newcommand{\bY}{\mathbf{Y}}
\newcommand{\bZero}{\boldsymbol{0}}
\newcommand{\skewt}{skew-$t$\xspace}
\newcommand{\Skewt}{Skew-$t$\xspace}
\newcommand{\eref}[1]{(\ref{#1})}
\newcommand{\sref}[1]{Section~\ref{#1}}

% ---- Local shorthands ----
\newcommand{\bphi}{\boldsymbol{\varphi}}
\newcommand{\ind}{\mathbf{1}}

% ---- Theorem environments ----
\theoremstyle{plain}
\newtheorem{proposition}{Proposition}[section]
\newtheorem{theorem}[proposition]{Theorem}
\newtheorem{lemma}[proposition]{Lemma}
\newtheorem{corollary}[proposition]{Corollary}

\theoremstyle{definition}
\newtheorem{definition}[proposition]{Definition}

\theoremstyle{remark}
\newtheorem{remark}[proposition]{Remark}

\lstset{
    basicstyle=\ttfamily\small,
    keywordstyle=\color{blue},
    commentstyle=\color{gray}\itshape,
    stringstyle=\color{teal},
    showstringspaces=false,
    breaklines=true,
    frame=single,
    framesep=4pt,
    language=R
}

\title{Prior Support Mismatch for $\varphi_1$ under the AR(2)
       Stationarity Triangle:\\ Quantification and a Nested
       Reparameterisation}
\author{Yi-Xuan Xie}
\date{2026-07-13}

\begin{document}
\maketitle

\begin{abstract}
\noindent
The AR(2) extension in \texttt{update\_params\_ar2.R} reuses the AR(1) prior
scale $N(0, 0.5^2)$ for both $\varphi_1$ and $\varphi_2$, although the
marginal support of $\varphi_1$ doubles from $(-1, 1)$ to $(-2, 2)$ under the
stationarity triangle. This note quantifies the resulting prior penalty and
posterior shrinkage, shows the effect is negligible for every current
simulation setting, and records the partial-autocorrelation
reparameterisation that would make the AR(1) and AR(2) priors nested if a
high-persistence regime is ever added.
\end{abstract}

\section{Stationarity regions and marginal supports}\label{sec:regions}

\begin{itemize}
  \item The unit-variance latent AR($p$) prior restricts $\bphi$ to the
        stationarity region; for $p = 1$ and $p = 2$ these are
        \begin{equation}\label{eq:regions}
          S_1 = (-1, 1),
          \qquad
          S_2 = \bigl\{ (\varphi_1, \varphi_2) :
            \varphi_1 + \varphi_2 < 1,\;
            \varphi_2 - \varphi_1 < 1,\;
            |\varphi_2| < 1 \bigr\},
        \end{equation}
        with $S_2$ the open triangle with vertices $(-2, -1)$, $(2, -1)$,
        $(0, 1)$ (\texttt{check\_ar2\_stability}).
  \item The marginal supports under $S_2$ are therefore
        \begin{equation}\label{eq:supports}
          \varphi_1 \in (-2, 2),
          \qquad
          \varphi_2 \in (-1, 1),
        \end{equation}
        so the support of $\varphi_1$ is twice that of the AR(1) coefficient
        while $\varphi_2$ matches it.
\end{itemize}

\section{Prior specification: original versus port}\label{sec:priors}

\begin{itemize}
  \item The original AR(1) code (\texttt{biometrics/update\_params.R},
        \texttt{updatePhiTS}) uses a truncated-normal proposal on $S_1$ with
        Hastings correction and the prior
        \begin{equation}\label{eq:prior1}
          p_1(\varphi)
          \;\propto\;
          N\!\bigl( \varphi ;\, 0,\, 0.5^2 \bigr)\,
          \ind\{ \varphi \in S_1 \},
          \qquad
          \Phi\!\left( \tfrac{1}{0.5} \right)
          - \Phi\!\left( -\tfrac{1}{0.5} \right)
          \approx 0.954,
        \end{equation}
        i.e.\ about $95\%$ of the untruncated mass lies inside the support
        --- scale and support are matched.
  \item The AR(2) port (\texttt{updatePhiAR2TS}) places the same scale on
        both coordinates,
        \begin{equation}\label{eq:prior2}
          p_2(\bphi)
          \;\propto\;
          N\!\bigl( \varphi_1 ;\, 0,\, 0.5^2 \bigr)\,
          N\!\bigl( \varphi_2 ;\, 0,\, 0.5^2 \bigr)\,
          \ind\{ \bphi \in S_2 \},
        \end{equation}
        the truncation being realised implicitly by the silent-rejection
        proposal (Proposition~3.3 of \texttt{where\_we\_are.tex}).
  \item The arguments \texttt{prior.mean}, \texttt{prior.sd} of
        \texttt{updatePhiAR2TS} are never passed by any production caller
        (\texttt{updateTauTS\_AR2}, \texttt{updateZTS\_AR2},
        \texttt{updateKnotsTS\_AR2}); only \texttt{tests/} overrides them, so
        \eref{eq:prior2} is effectively hardcoded, as \eref{eq:prior1} is in
        the original.
  \item For $\varphi_2$ the pair (support, scale) is identical to the AR(1)
        case, so the calibration carries over; the mismatch is confined to
        $\varphi_1$.
\end{itemize}

\section{Quantifying the penalty on the high-persistence corner}
\label{sec:penalty}

\begin{itemize}
  \item Under \eref{eq:prior2} the log-prior deficit of a candidate
        $\varphi_1$ relative to $\varphi_1 = 0$ is
        \begin{equation}\label{eq:penalty}
          \log \frac{p_2(0, \varphi_2)}{p_2(\varphi_1, \varphi_2)}
          = \frac{\varphi_1^2}{2 \cdot 0.5^2}
          = 2\, \varphi_1^2,
        \end{equation}
        e.g.\ $2.0$ nats at $\varphi_1 = 1$, $4.5$ nats at $\varphi_1 = 1.5$
        ($3\sigma$), and $8.0$ nats at the vertex-adjacent $\varphi_1 \to 2$.
  \item The region so penalised is the real-root, high-persistence corner of
        $S_2$ (both roots of $1 - \varphi_1 u - \varphi_2 u^2$ near the unit
        circle, slowly decaying monotone ACF), which is thus excluded a
        priori rather than by the data.
  \item All data-generating values in \texttt{simstudy/setup.R} satisfy
        $|\varphi_1| \leq 0.80 = 1.6\sigma$, i.e.\
        $2 \varphi_1^2 \leq 1.28$ nats, well inside the prior's effective
        support:
        \begin{equation}\label{eq:dgps}
          \bphi \in \bigl\{ (0.80, -0.35),\; (0.12, -0.05),\;
                            (0.15, 0.80) \bigr\}.
        \end{equation}
\end{itemize}

\section{Measured posterior shrinkage}\label{sec:shrinkage}

\begin{itemize}
  \item To isolate the prior's contribution from finite-sample bias, the same
        componentwise MH sampler (conditional likelihood, silent rejection,
        $T = 50$, $30$ replicates) was run under \eref{eq:prior2} and under a
        near-flat reference $N(0, 10^2)$ restricted to $S_2$;
        \cref{tab:shrink} reports posterior-mean averages.
\end{itemize}

\begin{table}[ht]
\centering
\caption{Posterior means of $(\varphi_1, \varphi_2)$ under the production
         prior \eref{eq:prior2} and a near-flat prior on $S_2$; $T = 50$,
         $30$ replicates, $K$ independent series per dataset.}
\label{tab:shrink}
\begin{tabular}{llccc}
\toprule
Setting & $K$ & Truth & Prior $\mathrm{sd} = 0.5$ & Flat
  ($\mathrm{sd} = 10$) \\
\midrule
16 near-unit-root & 1 & $(0.15,\; 0.80)$  & $(0.17,\; 0.73)$  & $(0.16,\; 0.74)$ \\
16 near-unit-root & 5 & $(0.15,\; 0.80)$  & $(0.16,\; 0.79)$  & $(0.15,\; 0.79)$ \\
9 strong          & 1 & $(0.80,\, -0.35)$ & $(0.73,\, -0.31)$ & $(0.77,\, -0.35)$ \\
9 strong          & 5 & $(0.80,\, -0.35)$ & $(0.78,\, -0.35)$ & $(0.79,\, -0.36)$ \\
10 weak           & 1 & $(0.12,\, -0.05)$ & $(0.10,\, -0.02)$ & $(0.10,\, -0.02)$ \\
10 weak           & 5 & $(0.12,\, -0.05)$ & $(0.12,\, -0.06)$ & $(0.12,\, -0.06)$ \\
\bottomrule
\end{tabular}
\end{table}

\begin{remark}[Interpretation]\label{rem:shrink}
The prior-attributable shrinkage peaks at $0.04$ (setting~9, $K = 1$,
$\varphi_1$: $0.73$ versus $0.77$) and vanishes at $K = 5$, where the
likelihood carries five times the information. The residual gap from truth
under the flat prior (e.g.\ $0.77$ versus $0.80$) is finite-sample AR bias,
which no prior choice removes; the current specification is benign for every
setting in \eref{eq:dgps}.
\end{remark}

\section{A caveat for the AR(1)-versus-AR(2) comparison}\label{sec:caveat}

\begin{itemize}
  \item The study's central comparison fits both an AR(1) and an AR(2)
        analysis model, whose coefficient priors are \eref{eq:prior1} and
        \eref{eq:prior2} respectively.
  \item These are not two versions of one prior: \eref{eq:prior2} restricted
        to $\{\varphi_2 = 0\}$ has $\varphi_1$-density proportional to
        $N(0, 0.5^2)$ on $(-1, 1)$, but the two models' implied priors on the
        ACF $\{\rho(h)\}_{h \geq 1}$ are not linked by any conditioning or
        marginalisation, so order comparisons carry a (small) prior confound.
  \item The confound is bounded by the shrinkage in \cref{tab:shrink}
        ($\leq 0.04$ on posterior means) and is therefore immaterial at the
        current design, but it should be stated when the comparison is
        written up.
\end{itemize}

\section{Nested priors via partial autocorrelations}\label{sec:pacf}

\begin{itemize}
  \item The clean resolution, if a high-persistence regime
        ($\varphi_1 > 1$) is ever added, is the partial-autocorrelation
        (PACF) reparameterisation of Barndorff-Nielsen and Schou (1973) and
        Monahan (1984): for $p = 2$,
        \begin{equation}\label{eq:pacf}
          \varphi_2 = r_2,
          \qquad
          \varphi_1 = r_1 (1 - r_2),
          \qquad
          (r_1, r_2) \in (-1, 1)^2,
        \end{equation}
        a smooth bijection from $(-1, 1)^2$ onto $S_2$ whose inverse is
        $r_1 = \varphi_1 / (1 - \varphi_2)$, $r_2 = \varphi_2$.
  \item Note $r_1 = \varphi_1 / (1 - \varphi_2) = \gamma_1$, the lag-one
        autocorrelation of the process, so the parameters are directly
        interpretable and the map is available in closed form in both
        directions.
  \item Placing independent priors $r_k \sim N(0, 0.5^2)$ truncated to
        $(-1, 1)$ yields a prior supported exactly on $S_2$ with no silent
        mass loss, and its geometry respects the triangle automatically:
        \begin{equation}\label{eq:nested}
          p = 1: \quad r_1 = \varphi \sim \eref{eq:prior1},
          \qquad
          p = 2: \quad (r_1, r_2) \ \text{i.i.d.\ as above},
        \end{equation}
        so the AR(1) prior is recovered exactly on the section $r_2 = 0$ ---
        the priors become nested and comparable across orders.
  \item Implementation cost is small (transform in
        \texttt{updatePhiAR2TS}, Jacobian
        $|\partial \bphi / \partial \mathbf{r}| = 1 - r_2$ in the acceptance
        ratio), but per \cref{rem:shrink} \emph{no code change is needed
        now}; this section exists so the option is on record.
\end{itemize}

\section{Status and decision}\label{sec:status}

\begin{itemize}
  \item \textbf{Decision (2026-07-15): no code change.} Unlike the
        $\bphi$-block likelihood issue (fixed the same day; see
        \texttt{tex/ar2\_phi\_exact\_likelihood/}), this is a prior
        \emph{elicitation} judgement, not a correctness defect: the sampler
        validly targets the posterior induced by \eref{eq:prior2}.
  \item The measured impact bound (\cref{tab:shrink}, $\leq 0.04$ and
        vanishing at $K = 5$) covers every DGP value in \eref{eq:dgps}, so
        the specification stays as is for the current study.
  \item Two follow-ups are on record: state the prior non-nestedness caveat
        (\sref{sec:caveat}) in the paper's AR(1)-versus-AR(2) comparison,
        and revisit the PACF reparameterisation (\sref{sec:pacf}) only if a
        DGP with $|\varphi_1| > 1$ is added.
\end{itemize}

\begin{thebibliography}{9}

\bibitem{bns1973}
Barndorff-Nielsen, O.\ and Schou, G.\ (1973).
On the parametrization of autoregressive models by partial autocorrelations.
\emph{Journal of Multivariate Analysis}, 3(4), 408--419.

\bibitem{monahan1984}
Monahan, J.~F.\ (1984).
A note on enforcing stationarity in autoregressive-moving average models.
\emph{Biometrika}, 71(2), 403--404.

\end{thebibliography}

\end{document}
